% Section 2 — The Three Times (compressed v3)

\section{The Three Times}
\label{sec:three-times}

We decompose any agent trajectory into three ontologically distinct projections of time (Figure~\ref{fig:three-times}).

\begin{figure}[h]
\centering
\includegraphics[width=\linewidth]{figures/three_times.pdf}
\caption{The Three Times ontology. $\twall$ (top, blue) is external and continuous. $\tstep$ (middle, red) is discrete policy invocations $a_0, a_1, \ldots$. $\tself$ (bottom, orange) is the agent's narrative of its own duration. A chronoceptively grounded policy enforces $\twall \approx \tstep \cdot \langle\Delta t\rangle \approx \tself$; the Augustine Problem (dashed arrow) is the failure of this identity. Each axis carries its named law (L1, L2, L3).}
\label{fig:three-times}
\end{figure}

\begin{definition}[Three Times]
\label{def:three-times}
For agent $A$ executing trajectory $\trajectory = (s_0, a_0, s_1, a_1, \ldots, s_n)$: $\twall(\trajectory) := t(s_n) - t(s_0)$ (external clock); $\tstep(\trajectory) := n$ (policy invocations); $\tself(\trajectory) := \selfparser(s_n^{\text{output}})$ (agent's reported duration; parser specified in Appendix~\ref{app:annotation-protocol}).
\end{definition}

Grounded chronoception requires $\twall \approx \tstep \cdot \langle\Delta t\rangle \approx \tself$, where $\langle\Delta t\rangle$ is the average per-step wall-clock cost. Augustine-problematic agents (\S\ref{sec:augustine-problem}) are those whose policies fail to enforce this identity --- a measurable property on single-turn trajectories. We use \emph{chronoception} in the cognitive-science sense (perception of one's own work duration), distinct from \citet{goel2025chronocept}'s use for fact validity and from temporal-reasoning benchmarks \citep{chu2024timebench} that test reasoning over external timelines.
